\documentclass[11pt]{article}
\usepackage[T1]{fontenc}
\usepackage[utf8]{inputenc}
\usepackage{amsmath,amssymb,amsthm}
\usepackage[margin=1in]{geometry}
\newtheorem{theorem}{Theorem}
\title{A Rational Two-Point Counterexample to the Zero-Radical Clause\\for TLMC 00000002617}
\author{Prepared for the account \texttt{idealistichacker}}
\date{}
\begin{document}
\maketitle

\section*{Frozen statement and scope}
The organizer source was frozen before this analysis:
\begin{flushleft}
\footnotesize
Source: \texttt{conjectures/00000002617.md}\\
Commit: \texttt{95acb520ec5607c826b8a997b1ef2fc82d6f7c57}\\
Git blob: \texttt{e23e80d92bfbf3e0970c427716f7ce3c2f97a39e}\\
SHA-256: \texttt{03ba0f9febb435b53de0c4df73642bbb647b1a0b9eef3f297cedf5263564bc5b}
\end{flushleft}

We address only the zero-radical clause at the conventional, fixed instance
\[
 A=\operatorname{IncidenceAlgebra}(\mathbb{Q},\operatorname{Fin}\,2),
 \qquad \operatorname{Fin}\,2=\{0<1\}.
\]
This is a finite locally finite chain with rational coefficients. We do not
make a statement about arbitrary coefficient rings or arbitrary posets, and
neither define nor use a notion of semisimplicity. The source's Möbius-function
language is not a proof object in this package.

\section*{Calculation in the actual incidence algebra}
Let \(e\in A\) take the value \(1\) on the strict interval \((0,1)\)
and \(0\) elsewhere. This is an actual Mathlib
\texttt{IncidenceAlgebra} function, not a substitute coordinate algebra.
The Lean file \texttt{lean4/Main.lean} expands Mathlib's
\texttt{IncidenceAlgebra.mul\_apply} over all pairs in
\(\operatorname{Fin}\,2\) and proves
\[
 (ye)^2=0 \quad\text{for every }y\in A.
\]
It also proves \(e\ne0\).

\section*{Jacobson-radical bridge}
Mathlib's actual \texttt{Ideal.mem\_jacobson\_iff} characterizes
membership of \(x\) in the Jacobson radical of an ideal. For each \(y\),
choose \(z=1-ye\). Since \((ye)^2=0\), its required expression simplifies to
\[
 zye+z-1=(ye-(ye)^2)+(1-ye)-1=0.
\]
Thus Lean proves
\[
 e\in\operatorname{Ideal.jacobson}(\bot:\operatorname{Ideal}\,A).
\]
Since \(e\ne0\), it follows that
\[
 \operatorname{Ideal.jacobson}(\bot:\operatorname{Ideal}\,A)\ne\bot.
\]
Hence the zero-radical clause has a counterexample under the explicitly fixed
rational two-point interpretation.

\section*{Lean evidence and trust boundary}
The project pins Lean 4.33.1 and Mathlib Git commit
\texttt{0df444a360eaa60ab8c11dca51a86af692955474}. \texttt{Check.lean} prints the transitive axiom footprints
for all advertised theorems. The current proofs use the standard Lean/Mathlib
principles \texttt{propext}, \texttt{Classical.choice}, and
\texttt{Quot.sound}; no new axiom, \texttt{sorry}, \texttt{admit},
\texttt{native\_decide}, or unsafe evaluation is used.

\section*{Status}
This package includes a PDF compiled from this source with Tectonic~0.15.0 and
deterministic PDF timestamp metadata. It has no independent review record and
does not claim organizer approval, merge, acceptance, ranking, award, or human
verification. Before any future publication, a new independent review, fresh
upstream duplicate/source checks, and release-coordinator validation are required.
\end{document}
